% Options for packages loaded elsewhere
% Options for packages loaded elsewhere
\PassOptionsToPackage{unicode}{hyperref}
\PassOptionsToPackage{hyphens}{url}
\PassOptionsToPackage{dvipsnames,svgnames,x11names}{xcolor}
%
\documentclass[
  12pt,
  letterpaper,
  DIV=11,
  numbers=noendperiod]{scrreprt}
\usepackage{xcolor}
\usepackage[margin = 1in]{geometry}
\usepackage{amsmath,amssymb}
\setcounter{secnumdepth}{-\maxdimen} % remove section numbering
\usepackage{iftex}
\ifPDFTeX
  \usepackage[T1]{fontenc}
  \usepackage[utf8]{inputenc}
  \usepackage{textcomp} % provide euro and other symbols
\else % if luatex or xetex
  \usepackage{unicode-math} % this also loads fontspec
  \defaultfontfeatures{Scale=MatchLowercase}
  \defaultfontfeatures[\rmfamily]{Ligatures=TeX,Scale=1}
\fi
\usepackage[]{baskervald}
\ifPDFTeX\else
  % xetex/luatex font selection
\fi
% Use upquote if available, for straight quotes in verbatim environments
\IfFileExists{upquote.sty}{\usepackage{upquote}}{}
\IfFileExists{microtype.sty}{% use microtype if available
  \usepackage[]{microtype}
  \UseMicrotypeSet[protrusion]{basicmath} % disable protrusion for tt fonts
}{}
\makeatletter
\@ifundefined{KOMAClassName}{% if non-KOMA class
  \IfFileExists{parskip.sty}{%
    \usepackage{parskip}
  }{% else
    \setlength{\parindent}{0pt}
    \setlength{\parskip}{6pt plus 2pt minus 1pt}}
}{% if KOMA class
  \KOMAoptions{parskip=half}}
\makeatother
% Make \paragraph and \subparagraph free-standing
\makeatletter
\ifx\paragraph\undefined\else
  \let\oldparagraph\paragraph
  \renewcommand{\paragraph}{
    \@ifstar
      \xxxParagraphStar
      \xxxParagraphNoStar
  }
  \newcommand{\xxxParagraphStar}[1]{\oldparagraph*{#1}\mbox{}}
  \newcommand{\xxxParagraphNoStar}[1]{\oldparagraph{#1}\mbox{}}
\fi
\ifx\subparagraph\undefined\else
  \let\oldsubparagraph\subparagraph
  \renewcommand{\subparagraph}{
    \@ifstar
      \xxxSubParagraphStar
      \xxxSubParagraphNoStar
  }
  \newcommand{\xxxSubParagraphStar}[1]{\oldsubparagraph*{#1}\mbox{}}
  \newcommand{\xxxSubParagraphNoStar}[1]{\oldsubparagraph{#1}\mbox{}}
\fi
\makeatother


\usepackage{longtable,booktabs,array}
\usepackage{calc} % for calculating minipage widths
% Correct order of tables after \paragraph or \subparagraph
\usepackage{etoolbox}
\makeatletter
\patchcmd\longtable{\par}{\if@noskipsec\mbox{}\fi\par}{}{}
\makeatother
% Allow footnotes in longtable head/foot
\IfFileExists{footnotehyper.sty}{\usepackage{footnotehyper}}{\usepackage{footnote}}
\makesavenoteenv{longtable}
\usepackage{graphicx}
\makeatletter
\newsavebox\pandoc@box
\newcommand*\pandocbounded[1]{% scales image to fit in text height/width
  \sbox\pandoc@box{#1}%
  \Gscale@div\@tempa{\textheight}{\dimexpr\ht\pandoc@box+\dp\pandoc@box\relax}%
  \Gscale@div\@tempb{\linewidth}{\wd\pandoc@box}%
  \ifdim\@tempb\p@<\@tempa\p@\let\@tempa\@tempb\fi% select the smaller of both
  \ifdim\@tempa\p@<\p@\scalebox{\@tempa}{\usebox\pandoc@box}%
  \else\usebox{\pandoc@box}%
  \fi%
}
% Set default figure placement to htbp
\def\fps@figure{htbp}
\makeatother





\setlength{\emergencystretch}{3em} % prevent overfull lines

\providecommand{\tightlist}{%
  \setlength{\itemsep}{0pt}\setlength{\parskip}{0pt}}



 


\KOMAoption{captions}{tableheading}
\makeatletter
\@ifpackageloaded{tcolorbox}{}{\usepackage[skins,breakable]{tcolorbox}}
\@ifpackageloaded{fontawesome5}{}{\usepackage{fontawesome5}}
\definecolor{quarto-callout-color}{HTML}{909090}
\definecolor{quarto-callout-note-color}{HTML}{0758E5}
\definecolor{quarto-callout-important-color}{HTML}{CC1914}
\definecolor{quarto-callout-warning-color}{HTML}{EB9113}
\definecolor{quarto-callout-tip-color}{HTML}{00A047}
\definecolor{quarto-callout-caution-color}{HTML}{FC5300}
\definecolor{quarto-callout-color-frame}{HTML}{acacac}
\definecolor{quarto-callout-note-color-frame}{HTML}{4582ec}
\definecolor{quarto-callout-important-color-frame}{HTML}{d9534f}
\definecolor{quarto-callout-warning-color-frame}{HTML}{f0ad4e}
\definecolor{quarto-callout-tip-color-frame}{HTML}{02b875}
\definecolor{quarto-callout-caution-color-frame}{HTML}{fd7e14}
\makeatother
\makeatletter
\@ifpackageloaded{caption}{}{\usepackage{caption}}
\AtBeginDocument{%
\ifdefined\contentsname
  \renewcommand*\contentsname{Table of contents}
\else
  \newcommand\contentsname{Table of contents}
\fi
\ifdefined\listfigurename
  \renewcommand*\listfigurename{List of Figures}
\else
  \newcommand\listfigurename{List of Figures}
\fi
\ifdefined\listtablename
  \renewcommand*\listtablename{List of Tables}
\else
  \newcommand\listtablename{List of Tables}
\fi
\ifdefined\figurename
  \renewcommand*\figurename{Figure}
\else
  \newcommand\figurename{Figure}
\fi
\ifdefined\tablename
  \renewcommand*\tablename{Table}
\else
  \newcommand\tablename{Table}
\fi
}
\@ifpackageloaded{float}{}{\usepackage{float}}
\floatstyle{ruled}
\@ifundefined{c@chapter}{\newfloat{codelisting}{h}{lop}}{\newfloat{codelisting}{h}{lop}[chapter]}
\floatname{codelisting}{Listing}
\newcommand*\listoflistings{\listof{codelisting}{List of Listings}}
\makeatother
\makeatletter
\makeatother
\makeatletter
\@ifpackageloaded{caption}{}{\usepackage{caption}}
\@ifpackageloaded{subcaption}{}{\usepackage{subcaption}}
\makeatother
\usepackage{bookmark}
\IfFileExists{xurl.sty}{\usepackage{xurl}}{} % add URL line breaks if available
\urlstyle{same}
\hypersetup{
  pdftitle={BRUV: Guia rápido de campo},
  pdfauthor={Pesquisa Meros SC},
  colorlinks=true,
  linkcolor={blue},
  filecolor={Maroon},
  citecolor={Blue},
  urlcolor={Blue},
  pdfcreator={LaTeX via pandoc}}


\title{BRUV: Guia rápido de campo}
\usepackage{etoolbox}
\makeatletter
\providecommand{\subtitle}[1]{% add subtitle to \maketitle
  \apptocmd{\@title}{\par {\large #1 \par}}{}{}
}
\makeatother
\subtitle{Projeto Meros do Brasil}
\author{Pesquisa Meros SC}
\date{}
\begin{document}
\maketitle
\begin{abstract}
Leia sempre antes de ir a campo para revisar os procedimentos.
\end{abstract}


\chapter{Procedimentos Pré-Campo}\label{procedimentos-pruxe9-campo}

\begin{itemize}
\item
  Verificar carga das baterias.
\item
  Verificar armazenamento dos cartões de memória.
\item
  Verificar funcionamento das câmeras.
\item
  Verificar preseça de todos os componentes e peças do \emph{BRUV}.
\item
  Verificar presença de todas as porcas para ajuste das .
\end{itemize}

\chapter{Procedimentos de Campo}\label{procedimentos-de-campo}

\section{Preparação para o
lançamento}\label{preparauxe7uxe3o-para-o-lanuxe7amento}

\begin{itemize}
\item
  Monte os pés do BRUV.
\item
  Prepare a haste da isca com o cano de armazenamento.
\item
  \begin{tcolorbox}[enhanced jigsaw, breakable, rightrule=.15mm, colback=white, bottomrule=.15mm, opacitybacktitle=0.6, left=2mm, coltitle=black, leftrule=.75mm, toptitle=1mm, toprule=.15mm, bottomtitle=1mm, titlerule=0mm, title=\textcolor{quarto-callout-tip-color}{\faLightbulb}\hspace{0.5em}{Tip}, colframe=quarto-callout-tip-color-frame, arc=.35mm, colbacktitle=quarto-callout-tip-color!10!white, opacityback=0]

  \textbf{Não adicione a isca ainda, para não ficar pingando no barco.
  Deixe para pôr no último momento.}

  \end{tcolorbox}
\item
  Libere o cabo vermelho em comformidade com a profudidade local
  estimada. Ex: para 15m de profundidade, libere 17 braças ( 1 braça
  \textasciitilde1.5m).
\item
  Amare a \textbf{ponta solta} do cabo de lançamento (\textbf{corda
  vermelha}) na argola da bóia de meia água do BRUV.
\item
  Amarre a bóia de sinalização de superfície (\textbf{bóia laranja}) ao
  cabo de lançamento.
\item
  \begin{tcolorbox}[enhanced jigsaw, breakable, rightrule=.15mm, colback=white, bottomrule=.15mm, opacitybacktitle=0.6, left=2mm, coltitle=black, leftrule=.75mm, toptitle=1mm, toprule=.15mm, bottomtitle=1mm, titlerule=0mm, title=\textcolor{quarto-callout-caution-color}{\faFire}\hspace{0.5em}{Caution}, colframe=quarto-callout-caution-color-frame, arc=.35mm, colbacktitle=quarto-callout-caution-color!10!white, opacityback=0]

  \textbf{NÃO} passe o cabo \textbf{SOMENTE} no D-ring plástico da bóia
  de sinalização, Passe também por dentro do tirante horizontal da bóia
  como backup de segurança para caso do D-ring quebre.

  \end{tcolorbox}
\item
  Retire ambas as câmeras e ligue-as (\textbf{Pressione o botão
  lateral}).
\item
  Conecte a câmera ao celular \textbf{via bluetooth} (Ver Quick GoPro
  App.)
\item
  \textbf{Formate os cartões de memória}.
\item
  Ponha para gravar.
\item
  Verifique se de fato está gravando.
\item
  \begin{tcolorbox}[enhanced jigsaw, breakable, rightrule=.15mm, colback=white, bottomrule=.15mm, opacitybacktitle=0.6, left=2mm, coltitle=black, leftrule=.75mm, toptitle=1mm, toprule=.15mm, bottomtitle=1mm, titlerule=0mm, title=\textcolor{quarto-callout-tip-color}{\faLightbulb}\hspace{0.5em}{Tip}, colframe=quarto-callout-tip-color-frame, arc=.35mm, colbacktitle=quarto-callout-tip-color!10!white, opacityback=0]

  Veja o monitor. Quando está gravando, a GoPro emite um luz vermelha
  intermitente.

  \end{tcolorbox}
\item
  Devolva as câmaras aos seus RESPECTIVOS cases.
\item
  \textbf{Cuidado para não por de cabeça para baixo!}
\item
  \begin{tcolorbox}[enhanced jigsaw, breakable, rightrule=.15mm, colback=white, bottomrule=.15mm, opacitybacktitle=0.6, left=2mm, coltitle=black, leftrule=.75mm, toptitle=1mm, toprule=.15mm, bottomtitle=1mm, titlerule=0mm, title=\textcolor{quarto-callout-tip-color}{\faLightbulb}\hspace{0.5em}{Tip}, colframe=quarto-callout-tip-color-frame, arc=.35mm, colbacktitle=quarto-callout-tip-color!10!white, opacityback=0]

  Se atente as marcações no bruv, elas devem coincidir. Lembre-se
  \textbf{vermelho} é bombordo (esquerda) e \textbf{verde} é boreste
  (direita).

  \end{tcolorbox}
\item
  Adicione a isca ao cano e acople a haste da isca no BRUV.
\item
  \begin{tcolorbox}[enhanced jigsaw, breakable, rightrule=.15mm, colback=white, bottomrule=.15mm, opacitybacktitle=0.6, left=2mm, coltitle=black, leftrule=.75mm, toptitle=1mm, toprule=.15mm, bottomtitle=1mm, titlerule=0mm, title=\textcolor{quarto-callout-caution-color}{\faFire}\hspace{0.5em}{Caution}, colframe=quarto-callout-caution-color-frame, arc=.35mm, colbacktitle=quarto-callout-caution-color!10!white, opacityback=0]

  Quando for conectar a haste da isca, tenha cuidado com a bóia de meia
  água do BRUV. Está deve passar por fora da haste. Lembre-se de
  levantá-la antes de fixar a haste.

  \end{tcolorbox}
\item
  VERIFIQUE SE O PROFUDIMETRO ESTA MARCANDO 0m. Caso não esteja, retorne
  o ponteiro a marcação inicial.
\end{itemize}

\section{Lançamento na água}\label{lanuxe7amento-na-uxe1gua}

\begin{itemize}
\item
  Tudo pronto! Vamos lançar pra água!
\item
  Precisa de pelo menos duas pessoas, uma para lançar e outra para safar
  o cabo. Bom mesmo é três, porque fica um atento as anotações e coleta
  do ponto no GPS (protocolo a criar)
\item
  Apoie o BRUV no bordo, metade das pernas dentro, metade fora, com a
  isca virada para o mar.

  \begin{tcolorbox}[enhanced jigsaw, breakable, rightrule=.15mm, colback=white, bottomrule=.15mm, opacitybacktitle=0.6, left=2mm, coltitle=black, leftrule=.75mm, toptitle=1mm, toprule=.15mm, bottomtitle=1mm, titlerule=0mm, title=\textcolor{quarto-callout-warning-color}{\faExclamationTriangle}\hspace{0.5em}{Warning}, colframe=quarto-callout-warning-color-frame, arc=.35mm, colbacktitle=quarto-callout-warning-color!10!white, opacityback=0]

  Não solte na popa, ao menos que o motor da embarcação esteja
  desligado, pois corre risco de enroscar no hélice.

  \end{tcolorbox}
\item
  Em caso de embarcação em movimento, solte o mais vertical e rápido
  possível, evitando assim muita derivação.
\item
  Solte-o e a medida que ganha profundidade realiza alguns puxões para
  ganhar prumo.
\item
  Se puxar demais ou entrar na água em uma posição ruim, o sistema bruv
  ficara de ponta cabeçã ou tompará para frente.
\item
  Já aconteceu de o sistema cair em pé, mas de alguma maneira tobar para
  trás apos muitos minutos em posição correta. Acredito que tenha sido o
  cabo curto demais em relação a ondulação da maré, mas é preciso
  verificar mais.
\item
  Quando sentir que a sonda chegou ao fundo, lançe o cabo restante, com
  a bóia de sinalização.

  \begin{tcolorbox}[enhanced jigsaw, breakable, rightrule=.15mm, colback=white, bottomrule=.15mm, opacitybacktitle=0.6, left=2mm, coltitle=black, leftrule=.75mm, toptitle=1mm, toprule=.15mm, bottomtitle=1mm, titlerule=0mm, title=\textcolor{quarto-callout-caution-color}{\faFire}\hspace{0.5em}{Caution}, colframe=quarto-callout-caution-color-frame, arc=.35mm, colbacktitle=quarto-callout-caution-color!10!white, opacityback=0]

  Não esqueça de marcar o ponto no GPS! É para cada lançamento!

  \end{tcolorbox}
\end{itemize}

\section{Que informações anotar?}\label{que-informauxe7uxf5es-anotar}

\begin{tcolorbox}[enhanced jigsaw, breakable, rightrule=.15mm, colback=white, bottomrule=.15mm, opacitybacktitle=0.6, left=2mm, coltitle=black, leftrule=.75mm, toptitle=1mm, toprule=.15mm, bottomtitle=1mm, titlerule=0mm, title=\textcolor{quarto-callout-important-color}{\faExclamation}\hspace{0.5em}{Important}, colframe=quarto-callout-important-color-frame, arc=.35mm, colbacktitle=quarto-callout-important-color!10!white, opacityback=0]

Leve sua prancheta ou caderneta de campo!

\end{tcolorbox}

\begin{itemize}
\item
  Data
\item
  Local
\item
  Participantes
\item
  Hora de lançamento
\item
  Hora de retirada
\item
  Ponto GPS
\item
  Profundidade
\item
  ID do BRUV (B1, B2 ou B3)
\end{itemize}

\section{Remoção da água}\label{remouxe7uxe3o-da-uxe1gua}

\begin{itemize}
\item
  Enquanto um dos integrantes puxa o sistema BRUV da água, o outro deve
  ir recolhendo a corda do convez.
\item
  Uma vez embarcado o equipamento:

  \begin{itemize}
  \item
    Anote a hora de retirada;
  \item
    Retire a bóia de sinalização;
  \item
    Solte e recolha o cabo;
  \item
    Registre a profundidade e retorne o ponteiro do profundímetro para a
    marcação de zero;
  \item
    Retire a haste que sustenta o cano da isca e dispense o material na
    água;
  \item
    Passe uma água no cano da isca para retirar o excesso de peixe
    morto;
  \item
    Caso não vá mais utilizar os BRUVs, desmonte e guarde-os nas bolsas
    de transporte;
  \item
    Caso ainda vá utilizar, deixe-os montados, sem a haste da isca e de
    forma organizada, em algum local do convez de modo a permitir a
    circulação da tripulação.
  \end{itemize}
\end{itemize}

\begin{tcolorbox}[enhanced jigsaw, breakable, rightrule=.15mm, colback=white, bottomrule=.15mm, opacitybacktitle=0.6, left=2mm, coltitle=black, leftrule=.75mm, toptitle=1mm, toprule=.15mm, bottomtitle=1mm, titlerule=0mm, title=\textcolor{quarto-callout-caution-color}{\faFire}\hspace{0.5em}{Caution}, colframe=quarto-callout-caution-color-frame, arc=.35mm, colbacktitle=quarto-callout-caution-color!10!white, opacityback=0]

Ao guardar os BRUVs na bolsa de transporte, tenha especial cuidado com
as lentes, de modo a evitar arranhões, pois estes interferem na análise
das imagens.

\end{tcolorbox}

\chapter{Pós-Campo}\label{puxf3s-campo}

\section{Lavagem}\label{lavagem}

\begin{itemize}
\item
  Lave toda a estrutura dos bruvs com água doce.
\item
  Faça com que a água entre nos tubos dos pés.
\end{itemize}

\section{Cartões de memória \&
Baterias}\label{cartuxf5es-de-memuxf3ria-baterias}

\begin{itemize}
\item
  Retire a bateria da câmera, isso facilita o acesso ao cartão de
  memória.
\item
  Retire o cartão de memória da câmera.
\item
  \textbf{Copie} todos os arquivos do cartão para o \textbf{drive
  externo} do Meros SC (Ver Transferindo arquivos ao HD.
\end{itemize}

\begin{tcolorbox}[enhanced jigsaw, breakable, rightrule=.15mm, colback=white, bottomrule=.15mm, opacitybacktitle=0.6, left=2mm, coltitle=black, leftrule=.75mm, toptitle=1mm, toprule=.15mm, bottomtitle=1mm, titlerule=0mm, title=\textcolor{quarto-callout-tip-color}{\faLightbulb}\hspace{0.5em}{Tip}, colframe=quarto-callout-tip-color-frame, arc=.35mm, colbacktitle=quarto-callout-tip-color!10!white, opacityback=0]

\textbf{Faça o quanto antes! Não deixe para fazer depois, pois terá
dificuldade de se lembrar dos metadados de campo.}

\end{tcolorbox}

\begin{itemize}
\item
  \textbf{NÃO formate os cartões pelo computador.} Isso deve ser feito
  via celular (Ver: Quick GoPro App).
\item
  Devolva os cartões as suas respectivas câmeras (Ver Foto).

  \begin{tcolorbox}[enhanced jigsaw, breakable, rightrule=.15mm, colback=white, bottomrule=.15mm, opacitybacktitle=0.6, left=2mm, coltitle=black, leftrule=.75mm, toptitle=1mm, toprule=.15mm, bottomtitle=1mm, titlerule=0mm, title=\textcolor{quarto-callout-tip-color}{\faLightbulb}\hspace{0.5em}{Tip}, colframe=quarto-callout-tip-color-frame, arc=.35mm, colbacktitle=quarto-callout-tip-color!10!white, opacityback=0]

  A dica é fazer este processo individualmente para cada BRUV. Escolha
  um dos lados para começar (E/D), depois realize cada etapa desta seção
  na câmera escolhida e só após concluido o passo-a-passo, siga para o
  outro lado. Assim você não irá se perder e cometer erros no trato das
  informações.

  \end{tcolorbox}
\item
  Guarde a bateria no case de campo no compartimento ``A carregar''.

  \begin{tcolorbox}[enhanced jigsaw, breakable, rightrule=.15mm, colback=white, bottomrule=.15mm, opacitybacktitle=0.6, left=2mm, coltitle=black, leftrule=.75mm, toptitle=1mm, toprule=.15mm, bottomtitle=1mm, titlerule=0mm, title=\textcolor{quarto-callout-tip-color}{\faLightbulb}\hspace{0.5em}{Tip}, colframe=quarto-callout-tip-color-frame, arc=.35mm, colbacktitle=quarto-callout-tip-color!10!white, opacityback=0]

  \textbf{Facilite. Deixe as baterias já carregadas e prontas para o
  próximo campo.}

  \end{tcolorbox}
\end{itemize}




\end{document}
